\section{Method}\label{sec:method}
Pauli propagation~\cite{ppframework} and tensor contraction~\cite{itensor} are central to our exact approach in deriving the LER. 
In short, Pauli propagation expresses the observables as sums of Pauli strings and tracks
their transformations under the operations. 
Firstly, we use Pauli propagation to calculate the effects of the \textit{incoming} error through the heavily non-Clifford gates dominated double-checking circuits. 
Then we encode the gates, noise and post-selection conditions in a tensor network, contract said tensor network to obtain $A(p)$ and $B(p)$ from Eqs.~\eqref{eq:probabilities}, which are then used to compute the LER.

\diagram[H]{method/pipeline}
{Steps of the calculation. Pauli propagation gives the response of each double-checking stage to an incoming error, grouping keeps the tables small, and one tensor network joins the stages. The network is contracted in two arithmetics, giving exact fractions at fixed $p$ and exact series coefficients.}
{fig:pipeline}
Figure~\ref{fig:pipeline} summarises the steps. Section~\ref{sec:pp} carries the check observables backwards through a double-checking stage, which gives the response of that stage to each incoming Pauli error. Section~\ref{sec:grouping} groups the errors that share a response, so that each stage becomes a small table. Section~\ref{sec:tensor-contraction} writes the gates, noise and post-selection as tensors, contracts them with the incoming labels open to fill each table, and joins the stages into $A(p)$ and $B(p)$. The same network is evaluated twice: with a rational $p$ in modular arithmetic, which gives the exact fractions, and with polynomials in $x=p/(1-p)$ truncated at degree ten, which gives the series coefficients.

\subsection{Method overview}
The $d=5$ cultivation circuit has two non-Clifford checking stages separated by a purely Clifford (code growth) sub-circuit, see Fig.~\ref{fig:overview} for an overview of the $d=5$ cultivation circuit and its stages/sub-circuits. 
Fig.~\ref{fig:d3-incoming} shows one of the smaller double-checking sub-circuit and incoming Pauli error (see Appendix~\ref{app:circuits} for both double-checking circuits). 

\widediagram{overview}{Stages of the full $d=5$ cultivation circuit. Ancillary sub-systems to perform measurements are dropped for simplicity.}{fig:overview}

\diagram[H]{method/d3_incoming_column}
{The $d=3$ double-checking circuit and the earlier stages.
The dashed boundary marks this circuit's input and the purple $X_7$ specifies an `incoming error'
present.}
{fig:d3-incoming}

Apart from the single injection gate, the non-Clifford gates occur in the two double-checking stages, and everything between them is Clifford, although the Clifford growth stage still requires the exact average of Appendix~\ref{app:growth}.

Here, `incoming' means present at the entrance to the double-checking stage, before its first $T$/$T^\dagger$ layer, after injection and any earlier stages. For the $d=3$ stage the earlier stages are Clifford apart from the injected state, so their effect on the ideal encoded magic state is a Pauli error $E$ occurring with some probability. The same description holds for the $d=5$ stage, because within each syndrome block the $d=3$ check leaves the ideal state up to a logical Pauli, see Appendix~\ref{app:growth}. We therefore describe the input as a fixed $E$ applied to the ideal state. For example, $E=X_7$ places an $X$ error on data qubit 7, marked in Fig.~\ref{fig:d3-incoming}. The probability of $E$ comes from the earlier stages. The noise inside double checking is averaged separately.

In these double-checking sub-circuits, for each incoming Pauli error $E$, we calculate two numbers (probabilities):
\begin{itemize}
\item $F_A$: the probability of passing post-selection, i.e. all trivial detectors and the specified noiseless final code readout, and
\item $F_B$: the probability of passing post-selection with a logical error in the output state.
\end{itemize}
We call the pair $(F_A,F_B)$ the \emph{response} to that error $E$.
The incoming error is held fixed, while the noise within the double-checking stage is averaged. This separation is what makes the response reusable: one table serves every fault pattern of the earlier stages.

In the full $d=5$ circuit, we retain the probabilities of the
different error labels leaving the earlier $d=3$ double-checking stage.
We then include the noise and post-selection during code growth
to obtain the probabilities of errors entering the final $d=5$ double-checking stage.
We multiply each final $d=5$ double-checking response by the
probability that its incoming error occurs and no earlier
post-selection detector fires. Summing these contributions
gives $A(p)$ and $B(p)$.

We will illustrate Pauli propagation briefly before describing the full calculation next.

\subsection{Pauli propagation}\label{sec:pp}

Pauli propagation~\cite{ppframework} computes an expectation value by carrying the measured observable backwards through the circuit and evaluating it on the input state, so a readout probability becomes a weighted sum of Pauli expectation values. For cultivation, Clifford gates do not increase the number of terms, and the structure of the double-checking stages lets us merge equivalent terms and sum over the noise directly. Rare logical errors therefore never need to be sampled: the LER is computed exactly, or as a series.
{\color{black}For a ZX-calculus description using Pauli semiwebs~\cite{zxflow}, see the discussion and Appendix~\ref{app:semiwebs}.}

As an example, we start by using $\rho$ as the input state to a circuit, $U$ the circuit unitary and
$O$ the observable to be measured. For a unitary circuit,
\begin{equation}
 \begin{gathered}
 \langle O\rangle=\Tr(U\rho U^\dagger O)
 =\Tr(\rho\,U^\dagger O U),\\
 U^\dagger O U=\sum_P c_P P.
 \end{gathered}
 \label{eq:heisenberg}
\end{equation}
The sum runs over tensor products of Pauli operators: $I,X,Y,Z$. A Clifford gate transforms
a Pauli string to another, possibly with a $\pm$ sign. A non-Clifford $T=\text{diag}(1,e^{i\pi/4})$ gate can
split one string into two `paths':
\begin{equation}
 \begin{gathered}
 T^\dagger X T=\frac{X{\color{red}-}Y}{\sqrt2},\quad
 T^\dagger Y T=\frac{X{\color{blue}+}Y}{\sqrt2},\\
 T=\operatorname{diag}(1,e^{i\pi/4}) \ .
 \end{gathered}
 \label{eq:T}
\end{equation}
The {\color{blue}blue ($+$)} and {\color{red}red ($-$)} signs in
Eq.~\eqref{eq:T} belong to the coefficients multiplying the Pauli
operators. Both terms are kept with their signs with nothing sampled. Each propagation path contributes the product
of its coefficients. Contributions to the same Pauli string are added
with their signs, so they can cancel.

Fig.~\ref{fig:pauli} follows an observable backwards through an actual
portion of the $d=3$ double-checking circuit, a single CNOT with circuit-level noise. Starting at $X_{13}$ on the
right to the left. The $T$ gate forks two
terms $X_{13}$ and $Y_{13}$, with coefficients ${\color{blue}+1/\sqrt2}$ and
${\color{red}-1/\sqrt2}$ respecitvely. The CNOT leaves $X_{13}$ unchanged and turns $Y_{13}$
into $Z_{14}Y_{13}$.

\diagram[h]{d3_pauli_steps}
{A $d=3$ double-checking circuit portion and its backward Pauli propagation.
The arrows carry coefficients.}
{fig:pauli}

\begingroup\color{black}
\begin{subequations}\label{eq:depolarising-channels}
The one-qubit depolarising channel is
\begin{equation}
\begin{aligned}
\mathcal N_1(\rho)
&=(1-p)\rho\\
&\quad+\frac{p}{3}(X\rho X+Y\rho Y+Z\rho Z).
\end{aligned}
\label{eq:depolarising-one}
\end{equation}
The two-qubit depolarising channel is
\begin{equation}
\begin{aligned}
\mathcal N_2(\rho)
&=(1-p)\rho\\
&\quad+\frac{p}{15}
\sum_{\substack{P,Q\in\{I,X,Y,Z\}\\(P,Q)\neq(I,I)}}
(P\otimes Q)\rho(P\otimes Q).
\end{aligned}
\label{eq:depolarising-two}
\end{equation}
\end{subequations}
\par\endgroup
Their physical Pauli faults have positive
probabilities, which we sum exactly. For example, 
\begin{equation}
 \mathcal N_1^\dagger(X)
 =(1-p)X+\frac p3(X-X-X)
 =\lambda_1X,
\end{equation}
where $\lambda_1=1-4p/3$. The dagger here means the adjoint channel, which propagates
the observable backwards through the noise operators~\cite{ppframework}. Similarly,
$\lambda_2=1-16p/15$ for a non-identity two-qubit Pauli operator\footnote{Explicitly, for $P,Q\in\{I,X,Y,Z\}$ and $(P,Q)\neq(I,I)$,
$\mathcal N_2^\dagger(P\otimes Q) =\lambda_2(P\otimes Q)$, where $\lambda_2=1-\frac{16p}{15}$.}. Following all four steps in Fig.~\ref{fig:pauli} gives
\begin{equation}
 X_{13}\longmapsto\frac{\lambda_1\lambda_2}{\sqrt2}
       (X_{13}-Z_{14}Y_{13}).
 \label{eq:localexample}
\end{equation}
To obtain a scalar, one can evaluate the resultant Pauli terms on an input
state and add their weighted contributions. For a Pauli measurement readout with outcome parities $\pm 1$,
the mean outcome determines the two probabilities:
\begin{equation}
 \Pr(O=\pm1)=\frac{1\pm\langle O\rangle}{2}.
 \label{eq:pauli-probability}
\end{equation}
For illustration, one can use input $\ket0_{14}\ket+_{13}$;
$\langle X_{13}\rangle=1$ and $\langle Z_{14}Y_{13}\rangle=0$.
The output mean is therefore $\lambda_1\lambda_2/\sqrt2$.

\subsubsection{Pauli propagation on the double-checking circuit}

We now follow the observable through the whole $d=3$ check of Fig.~\ref{fig:d3-incoming} at $p=0$, to show how the 128 terms produced by the entry layer collapse to a single logical observable. This collapse is what keeps the response tables small. Consider the central $X$-basis measurement on qubit $7$. 
Propagating $X_7$ backwards through the four CNOT layers $\mathtt{C}_4,\mathtt{C}_3,\mathtt{C}_2,\mathtt{C}_1$ gives the strings in
Fig.~\ref{fig:d3-full-propagation}(a). Every step still returns one Pauli string, and the CNOTs spread its support without introducing a sum.

Write $\mathcal D=\{0,3,7,9,10,12,13\}$ for the data qubits.
At the entry layer $D^\dagger=\prod_{q\in\mathcal D}T_q^\dagger$,
backward propagation uses $TXT^\dagger=(X+Y)/\sqrt2$.
The four ancillary $X$ factors have expectation $+1$ on their prepared
$\ket{+}$ states. The operator left on the data is:
\begin{equation}
 R=2^{-7/2}\prod_{q\in\mathcal D}(X_q+Y_q) \ .
 \label{eq:v2-central-operator}
\end{equation}
This product contains 128 Pauli strings. On a codespace input, 112 of the 128 Pauli strings flip one or more code stabiliser
and hence have zero expectation on the ideal encoded state.
Of the remaining 16 Pauli strings, eight act as $X_L$ and eight as $-Y_L$.
Their coefficients combine, forming the logical observable: $M_L$,
\begin{equation}
 \Pi_0R\Pi_0=M_L\Pi_0,\qquad M_L=\frac{X_L-Y_L}{\sqrt2} \ .
 \label{eq:v2-central-logical}
\end{equation}
The ideal state is then:
$\ket{T_L^\dagger}=(\ket{0_L}+e^{-i\pi/4}\ket{1_L})/\sqrt2$,
which has $\langle M_L\rangle=1$\footnote{Pauli operators: $X_L=\prod_{q\in\mathcal D}X_q$, $Z_L=\prod_{q\in\mathcal D}Z_q$,
and $Y_L=iX_LZ_L$.}.

\diagram{method/d3_full_propagation}
{(a) The central $X_7$ propagated backwards, with each string
beside its node. (b) The entry layer and code projection reduce
128 terms to two logical terms.}
{fig:d3-full-propagation}

The incoming error changes the state on which the propagated operator
is evaluated. For a specified $E$, the central readout probability is
\begin{equation}
 \Pr(m_7=0\mid E)
 =\frac{1+\langle T_L^\dagger|E^\dagger R E|T_L^\dagger\rangle}{2}.
 \label{eq:v2-incoming-mean}
\end{equation}
Substituting $E=I$ gives probability 1, as expected. For $E=X_7$ the conjugation turns $(X_7+Y_7)$ in Eq.~\eqref{eq:v2-central-operator} into $(X_7-Y_7)$ and the probability drops to $1/2$. This is only the readout of the central measurement on qubit 7 in Fig.~\ref{fig:d3-incoming}, however: even when it passes, the six auxiliary measurements and the final code projection detect the syndrome of $X_7$, so with no further faults the check rejects and $F_A(X_7;0)=0$. For $p>0$ a second fault can cancel that syndrome and let $X_7$ pass, which is why every incoming $E$, rejected or not at $p=0$, has an entry $F_A(E;p)$, $F_B(E;p)$ in the table.

With noise, the same backward propagation is averaged over every fault pattern inside the check, keeping only the measurement records with trivial detector parities (which means post-selection is successful). Call the resulting
unnormalised output state $\sigma_E$. Its trace equals the probability of obtaining measurement records
with trivial detector parities, for a fixed incoming error $E$ in either of the two double-checking subcircuits.
The noiseless final readout uses the projector $\Pi_0$ onto the
accepted code space and $\Pi_{\rm bad}$ onto the states within
that code space with the wrong logical outcome. A projector's expectation is the probability
of its event. Thus
\begin{equation}
 \begin{aligned}
 F_A(E;p)&=\Tr(\Pi_0\sigma_E),\\
 F_B(E;p)&=\Tr(\Pi_{\rm bad}\sigma_E).
 \end{aligned}
 \label{eq:response-projectors}
\end{equation}
We evaluate these traces by propagating the projectors backwards. The state $\sigma_E$ need not be stored.
Their Pauli expansions can have negative coefficients, but the final
traces are probabilities. See Appendix~\ref{app:circuits} for these projectors.

So far the input carried the error and the check was noiseless. We now reverse the roles: an ideal input, and a fixed pattern of faults inside the check, $X_3$ in the `fold' layers and $Y_9$ in the `unfold' layer (Fig.~\ref{fig:d3-fixed-faults}), chosen because it is the smallest pattern that passes post-selection with a logical error.
\diagram{method/d3_accepted_faults}
{A fixed two-fault pattern in the $d=3$ double-checking circuit.
Coloured boxes are $X_3$ and $Y_9$ errors with their propagation traced.}
{fig:d3-fixed-faults}
The first fault propagates to $X_0X_1X_3$, its ancillary $X_1$ does not
flip the final $X$ readout. The data error before the exit layer is
$Q=X_0X_3Y_9$. Forward propagation through that layer gives
\begin{equation}
 DQD^\dagger=\frac{(X_0+Y_0)(X_3+Y_3)(Y_9-X_9)}{\sqrt8}.
 \label{eq:v2-fixed-fault}
\end{equation}
Fig.~\ref{fig:d3-fault-expansion} shows all eight terms.
The code projection removes six. The surviving operator is
$-(X_L+Y_L)/\sqrt8=-N_L/2$, where
$N_L=(X_L+Y_L)/\sqrt2$ anticommutes with $M_L$.
The accepted state is therefore logically wrong: two faults lead to an accepted logical error, hence the fault distance is $d_\text{fault}=2$. 
A ZX-calculus contraction of the noiseless circuit with these two faults inserted~\cite{wan} confirms it. No detector fires with probability $1/4$, and every accepted output carries the logical error.

\widediagram{method/d3_fault_expansion}
{All eight signed terms for a fixed two-fault example.
Every edge contributes $1/\sqrt2$ with its indicated sign. The two surviving terms give acceptance $1/4$ and a logical error.}
{fig:d3-fault-expansion}


\subsection{Grouping equivalent Pauli errors}\label{sec:grouping}
\definecolor{tnblue}{RGB}{0,118,186}
\definecolor{tnorange}{RGB}{181,23,0}
\definecolor{tngreen}{RGB}{1,113,0}
\tikzset{
  tn tensor/.style={rectangle,rounded corners=2pt,draw=black,
    fill=tnblue,text=white,line width=1.1pt,align=center,inner sep=2mm,
    minimum height=9mm,font=\small},
  tn endpoint/.style={tn tensor},
  tn noise/.style={tn tensor,fill=tnorange},
  tn intermediate/.style={tn tensor,fill=tngreen},
  tn constraint/.style={tn tensor,fill=black!16,text=black},
  tn leg/.style={line width=.6pt,draw=black},
  tn bundle/.style={tn leg,double,double distance=.6pt},
  tn copy/.style={circle,fill=black,draw=none,inner sep=0pt,minimum size=2.4pt},
  tn operation/.style={-{Stealth[length=1.8mm]},line width=.6pt,draw=black!65},
  tn label/.style={font=\scriptsize,align=center,inner sep=.8mm},
  tn heading/.style={font=\small,align=center}
}


The injected and then cultivated magic state is encoded in a triangular colour code.
\begin{figure}%
\centering
\begin{minipage}[c]{.38\linewidth}\centering
\begin{adjustbox}{max width=\linewidth}\begin{tikzpicture}[x=1cm,y=1cm,line width=.6pt]
\definecolor{dthreeblue}{RGB}{58,80,190}
\definecolor{dthreered}{RGB}{219,51,51}
\definecolor{dthreegreen}{RGB}{51,166,64}
\coordinate (q0) at (0,3.6);
\coordinate (q3) at (-1.0392,1.8);
\coordinate (q7) at (0,1.2);
\coordinate (q9) at (-2.0784,0);
\coordinate (q10) at (1.0392,1.8);
\coordinate (q12) at (0,0);
\coordinate (q13) at (2.0784,0);
\filldraw[fill=dthreeblue!48] (q0)--(q3)--(q7)--(q10)--cycle;
\filldraw[fill=dthreered!48] (q3)--(q9)--(q12)--(q7)--cycle;
\filldraw[fill=dthreegreen!48] (q7)--(q12)--(q13)--(q10)--cycle;
\node[note] at (0,2.25) {$f_1$};
\node[note] at (-.92,.70) {$f_2$};
\node[note] at (.92,.70) {$f_3$};
\foreach \q in {0,3,7,9,10,12,13}
  \node[circle,draw,fill=white,inner sep=0pt,minimum size=5mm,
        font=\footnotesize] at (q\q) {$\q$};
\end{tikzpicture}
\end{adjustbox}
\par\smallskip $d=3$: seven data sites
\end{minipage}\hfill
\begin{minipage}[c]{.58\linewidth}\centering
\begin{adjustbox}{max width=\linewidth}\begin{tikzpicture}[x=.5cm,y=.866025403784cm]
\definecolor{codered}{HTML}{DB3333}
\definecolor{codegreen}{HTML}{33A640}
\definecolor{codeblue}{HTML}{3A50BE}
\coordinate (d5q0) at (0,0);
\coordinate (d5q3) at (1,1);
\coordinate (d5q7) at (3,1);
\coordinate (d5q9) at (3,3);
\coordinate (d5q11) at (4,0);
\coordinate (d5q13) at (4,2);
\coordinate (d5q15) at (4,4);
\coordinate (d5q17) at (6,0);
\coordinate (d5q19) at (6,2);
\coordinate (d5q21) at (6,4);
\coordinate (d5q23) at (6,6);
\coordinate (d5q25) at (7,1);
\coordinate (d5q27) at (7,3);
\coordinate (d5q29) at (7,5);
\coordinate (d5q32) at (9,1);
\coordinate (d5q34) at (9,3);
\coordinate (d5q36) at (10,0);
\coordinate (d5q38) at (10,2);
\coordinate (d5q40) at (12,0);
\filldraw[fill=codeblue!35,draw=codeblue!75!black,line width=.55pt] (d5q0) -- (d5q11) -- (d5q7) -- (d5q3) -- cycle;
\filldraw[fill=codeblue!35,draw=codeblue!75!black,line width=.55pt] (d5q13) -- (d5q19) -- (d5q27) -- (d5q21) -- (d5q15) -- (d5q9) -- cycle;
\filldraw[fill=codeblue!35,draw=codeblue!75!black,line width=.55pt] (d5q17) -- (d5q36) -- (d5q32) -- (d5q25) -- cycle;
\filldraw[fill=codered!35,draw=codered!75!black,line width=.55pt] (d5q3) -- (d5q7) -- (d5q13) -- (d5q9) -- cycle;
\filldraw[fill=codered!35,draw=codered!75!black,line width=.55pt] (d5q15) -- (d5q21) -- (d5q29) -- (d5q23) -- cycle;
\filldraw[fill=codered!35,draw=codered!75!black,line width=.55pt] (d5q25) -- (d5q32) -- (d5q38) -- (d5q34) -- (d5q27) -- (d5q19) -- cycle;
\filldraw[fill=codegreen!35,draw=codegreen!75!black,line width=.55pt] (d5q11) -- (d5q17) -- (d5q25) -- (d5q19) -- (d5q13) -- (d5q7) -- cycle;
\filldraw[fill=codegreen!35,draw=codegreen!75!black,line width=.55pt] (d5q27) -- (d5q34) -- (d5q29) -- (d5q21) -- cycle;
\filldraw[fill=codegreen!35,draw=codegreen!75!black,line width=.55pt] (d5q36) -- (d5q40) -- (d5q38) -- (d5q32) -- cycle;
\node[site,minimum size=6mm,inner sep=0pt,font=\scriptsize] at (d5q0) {$T$};
\node[font=\scriptsize,inner sep=0pt,outer sep=0pt,below=3.2mm] at (d5q0) {0};
\node[site,minimum size=6mm,inner sep=0pt,font=\scriptsize] at (d5q3) {$T^\dagger$};
\node[font=\scriptsize,inner sep=0pt,outer sep=0pt,below=3.2mm] at (d5q3) {3};
\node[site,minimum size=6mm,inner sep=0pt,font=\scriptsize] at (d5q7) {$T^\dagger$};
\node[font=\scriptsize,inner sep=0pt,outer sep=0pt,below=3.2mm] at (d5q7) {7};
\node[site,minimum size=6mm,inner sep=0pt,font=\scriptsize] at (d5q9) {$T^\dagger$};
\node[font=\scriptsize,inner sep=0pt,outer sep=0pt,below=3.2mm] at (d5q9) {9};
\node[site,minimum size=6mm,inner sep=0pt,font=\scriptsize] at (d5q11) {$T$};
\node[font=\scriptsize,inner sep=0pt,outer sep=0pt,below=3.2mm] at (d5q11) {11};
\node[site,minimum size=6mm,inner sep=0pt,font=\scriptsize] at (d5q13) {$T^\dagger$};
\node[font=\scriptsize,inner sep=0pt,outer sep=0pt,below=3.2mm] at (d5q13) {13};
\node[site,minimum size=6mm,inner sep=0pt,font=\scriptsize] at (d5q15) {$T^\dagger$};
\node[font=\scriptsize,inner sep=0pt,outer sep=0pt,below=3.2mm] at (d5q15) {15};
\node[site,minimum size=6mm,inner sep=0pt,font=\scriptsize] at (d5q17) {$T$};
\node[font=\scriptsize,inner sep=0pt,outer sep=0pt,below=3.2mm] at (d5q17) {17};
\node[site,minimum size=6mm,inner sep=0pt,font=\scriptsize] at (d5q19) {$T$};
\node[font=\scriptsize,inner sep=0pt,outer sep=0pt,below=3.2mm] at (d5q19) {19};
\node[site,minimum size=6mm,inner sep=0pt,font=\scriptsize] at (d5q21) {$T^\dagger$};
\node[font=\scriptsize,inner sep=0pt,outer sep=0pt,below=3.2mm] at (d5q21) {21};
\node[site,minimum size=6mm,inner sep=0pt,font=\scriptsize] at (d5q23) {$T^\dagger$};
\node[font=\scriptsize,inner sep=0pt,outer sep=0pt,below=3.2mm] at (d5q23) {23};
\node[site,minimum size=6mm,inner sep=0pt,font=\scriptsize] at (d5q25) {$T^\dagger$};
\node[font=\scriptsize,inner sep=0pt,outer sep=0pt,below=3.2mm] at (d5q25) {25};
\node[site,minimum size=6mm,inner sep=0pt,font=\scriptsize] at (d5q27) {$T^\dagger$};
\node[font=\scriptsize,inner sep=0pt,outer sep=0pt,below=3.2mm] at (d5q27) {27};
\node[site,minimum size=6mm,inner sep=0pt,font=\scriptsize] at (d5q29) {$T^\dagger$};
\node[font=\scriptsize,inner sep=0pt,outer sep=0pt,below=3.2mm] at (d5q29) {29};
\node[site,minimum size=6mm,inner sep=0pt,font=\scriptsize] at (d5q32) {$T^\dagger$};
\node[font=\scriptsize,inner sep=0pt,outer sep=0pt,above=3.2mm] at (d5q32) {32};
\node[site,minimum size=6mm,inner sep=0pt,font=\scriptsize] at (d5q34) {$T^\dagger$};
\node[font=\scriptsize,inner sep=0pt,outer sep=0pt,below=3.2mm] at (d5q34) {34};
\node[site,minimum size=6mm,inner sep=0pt,font=\scriptsize] at (d5q36) {$T$};
\node[font=\scriptsize,inner sep=0pt,outer sep=0pt,below=3.2mm] at (d5q36) {36};
\node[site,minimum size=6mm,inner sep=0pt,font=\scriptsize] at (d5q38) {$T^\dagger$};
\node[font=\scriptsize,inner sep=0pt,outer sep=0pt,below=3.2mm] at (d5q38) {38};
\node[site,minimum size=6mm,inner sep=0pt,font=\scriptsize] at (d5q40) {$T$};
\node[font=\scriptsize,inner sep=0pt,outer sep=0pt,below=3.2mm] at (d5q40) {40};
\end{tikzpicture}
\end{adjustbox}
\par\smallskip $d=5$: entry layer $D^\dagger$
\end{minipage}
\caption{Triangular colour codes, number are qubit indices from~\cite{clifft, soft}. On the right, each node shows its entry $T$ or $T^\dagger$ gate in one layer of gates in the $d=5$ double-checking stage.}
\label{fig:patches}
\end{figure}
The $[[7,1,3]]$ ($d=3$) variant code uses seven data qubits, while the $d=5$ cultivation circuit
grows this to a $[[19,1,5]]$ code. In Fig.~\ref{fig:patches}, each coloured face defines an
$X$ and a $Z$ stabiliser, formed by multiplying the corresponding
Pauli operators over the qubits on that face. The code space is the
simultaneous $+1$ eigenspace of these stabilisers~\cite{albert2026handbookerrorcorrectingcodes}. There are six
independent stabiliser generators at $d=3$ and eighteen at $d=5$,
with one logical qubit encoded in either case. Pauli errors that differ by a stabiliser (or product of stabilisers) have the same action
on the encoded magic state and share a response-table entry.

We need one response for each distinguishable incoming error, rather than
one for every Pauli string (or else we need to store exponentially many of them). 
To identify errors that share a response, we first represent each
incoming Pauli string by its $X$ and $Z$ components using two binary
masks, $\boldsymbol a$ and $\boldsymbol b$\footnote{At each qubit $q$,
the tuple $(a_q,b_q)$ maps to the Pauli operators:
$I\leftrightarrow(0,0)$,
$X\leftrightarrow(1,0)$,
$Z\leftrightarrow(0,1)$ and
$Y\leftrightarrow(1,1)$, ignoring the overall phase.}. 
The syndrome records which code stabilisers have been flipped to $-1$ parity (represented by bit $1$)\footnote{Parity $+1$ is represented by bit $0$.}.
For the $d=3$ patch in Fig.~\ref{fig:patches}, the parity-check matrix is
\begin{equation}
 H^{(3)}=
 \begin{pmatrix}
 1&1&1&0&1&0&0\\
 0&1&1&1&0&1&0\\
 0&0&1&0&1&1&1
 \end{pmatrix}.
\end{equation}
The columns follow qubit order $(0,\allowbreak3,\allowbreak7,\allowbreak9,\allowbreak10,\allowbreak12,\allowbreak13)$. The rows correspond to the blue, red and green faces.
Each row marks the qubits belonging to that face.

For either code, we label each incoming error by
\begin{equation}
 \begin{aligned}
 \alpha&=(H\boldsymbol a,\ \boldsymbol1\cdot\boldsymbol a),\\
 \beta&=(H\boldsymbol b,\ \boldsymbol1\cdot\boldsymbol b),
 \end{aligned}
 \label{eq:address}
\end{equation}
all in modulo two. Here $H\boldsymbol a$ and
$H\boldsymbol b$ are the $Z$- and $X$-stabiliser syndromes,
respectively. The final bit in each label is the parity of the
corresponding mask, with $\boldsymbol1$ being the all-ones vector, 
\[
\boldsymbol1 = (1,1,\ldots,1)^{\text{T}} \ .
\]
These two additional bits distinguish logical Pauli classes
with the same syndrome.
Errors with the same $(\alpha,\beta)$ differ by a stabiliser, which acts trivially on the encoded state, so they have the same response.

For example, $X_7$ and $X_0X_3X_{10}$ differ by the blue-face
stabiliser $X_0X_3X_7X_{10}$ and therefore have the same effect
on the encoded state.

\diagram{frame_example}{Two Pauli errors share a response-table entry at $d=3$.}{fig:frames}

At $d=3$, each label $\alpha$ or $\beta$ has four bits, so each response table has
$16^2=256$ entries. At $d=5$, each label $\alpha$ or $\beta$ has ten bits and a table has
$1024^2$ entries. Grouping stabiliser-equivalent errors reduces the number of
response-table entries from $4^7=16\,384$ to $256$ at $d=3$,
and from $4^{19}\approx2.75\times10^{11}$ to
$1\,048\,576$ at $d=5$. We calculate and store one response
for each equivalence class.
These error labels become tensor indices. Contracting the tensors combines
the double-checking responses with the noise from earlier stages to give
$A(p)$ and $B(p)$.

\subsection{Tensor contraction}\label{sec:tensor-contraction}
The response of a check to an incoming error is a sum over every fault configuration inside the check. Writing the gates, noise and post-selection conditions as small tensors organises that sum: contracting them performs the noise average, and leaving the incoming labels open computes all table entries at once, sharing the internal sums between them.
\diagram{tensor_definition}
{(a) A tensor with $r$ indices. (b) Contracting $j$ leaves all
$i_1,\ldots,i_m$ and $k_1,\ldots,k_n$ indices open.}
{fig:tensor-definition}
A tensor is an array whose entries are labelled by indices and defined as follows. For $r$ binary indices, the entries are written
$\mathcal T_{i_1\cdots i_r}$, with $i_s\in\{0,1\}$.
Each leg represents one index. Contracting tensors multiplies
their entries and sums over the connected indices. Open legs
label the indices that remain. A double line bundles several binary indices. Each leg is an index over error labels, not a physical qubit wire. Throughout this manuscript, we draw tensors as rounded boxes with legs as indices, contrasting with square white boxes elsewhere indicating calculation steps.

\subsubsection{Response tables as tensors}
For fixed $p$, the response table $F_A$ is a matrix where the row
index groups the bits of $\alpha$ and column index
groups the bits of $\beta$. Its entry gives the probability (for
that fixed incoming error) that the double-checking stage and noiseless final code readout both pass. The corresponding entry of $F_B$ also requires
a logical error. Both probabilities are averaged over the noise within the
double-checking stage.

Let $\mu(\alpha,\beta;p)$ be the joint probability that the earlier
stages pass post-selection and supply that incoming error. Then
\begin{equation}
 \begin{aligned}
 A(p)&=\sum_{\alpha,\beta}\mu(\alpha,\beta;p)F_A(\alpha,\beta;p),\\
 B(p)&=\sum_{\alpha,\beta}\mu(\alpha,\beta;p)F_B(\alpha,\beta;p).
 \end{aligned}
 \label{eq:responseaverage}
\end{equation}
For each pair $(\alpha,\beta)$, multiply its weight $\mu$
by the corresponding response $F_A$ or $F_B$, then sum over all pairs.
This is a tensor contraction, shown in Fig.~\ref{fig:response}(b).
Each weight $\mu$ is already the joint probability of passing the earlier post-selection and producing
the specified incoming error. Because the weights sum to the earlier acceptance probability, no further normalisation is needed before taking $\PL=B/A$, giving the LER.

\diagram{response}
{(a) Eight binary indices label the $d=3$ response table.
(b) Bundling them into $\alpha,\beta$ and summing with $\mu$ gives $A$ or $B$.
Fig.~\ref{fig:opennet} in Appendix~\ref{app:response}
shows the internal network.}
{fig:response}

\subsubsection{Tensors for the circuit operations}
The eight open legs in the tensor from Fig.~\ref{fig:response}(a) label a completed
response table. To calculate that table, we contract smaller tensors
for the gates, noise and post-selection, keeping the incoming error
labels open, so that the internal sums are shared between table entries. 
We choose the tensor contraction order carefully to keep intermediate tensors small
without changing the results (see Appendix~\ref{app:tensor-examples-v2} for details). Figure~\ref{fig:orderbubbles} shows one such order as nested bubbles and as a tree.
\diagram[H]{order_bubbles}{A contraction order for the noisy growth stage of Fig.~\ref{fig:overview}, drawn two ways. Its 569 factors, the $d=3$ table $W$, 567 grouped noise factors $g_i$ and the $d=5$ response $F_O$, share one parity constraint on their labels (Appendix~\ref{app:growth}). $\dots$ stand for the omitted factors and joins. (a) Nested bubbles: each bubble is joined into one intermediate table, stored on the legs that cross its boundary. (b) The same order as a tree, one node per join, the dotted node standing for joins 3 to 566, the root being the scalar. The order changes the table sizes, not the contracted result.}{fig:orderbubbles}

For example, the CNOT tensor $C$ in Fig.~\ref{fig:response-network}(a)
contains the two parity tensors expanded in panel (b).
An $x$ bit records the $X$ component
of an error. A $v$ bit is a Fourier index for its $Z$ component.
In Fourier form the CNOT acts on the $Z$ labels by the same parity rule as on the $X$ labels, and the noise becomes a multiplicative weight (Appendix~\ref{app:response}).
Each grey tensor is $1$ when its three bits have even parity, and $0$
otherwise. The vector $\boldsymbol u$ labels the error just after the
gate; $\boldsymbol x'$ labels it after the noise. The repeated Fourier
label $\boldsymbol v'$ is shared with the noise tensor $H_{2,p}$.
Noise supplies its weights as defined in Appendix~\ref{app:response}.
The two $T$ layers and final selection are compiled into signed
terminal factors.
Appendix~\ref{app:response} gives their equations and the final transform
from the internal Fourier indices to $F_A,F_B$.

\widediagram{method/cnot_factors}
{(a) A circuit portion and its CNOT/noise tensors.
(b) The faint box groups the two parity factors into $C$.
Black dots reuse the same index. Eq.~\eqref{eq:gate-noise-tensor} and Fig.~\ref{fig:opennet} in Appendix~\ref{app:response} place this pair inside the response network.}
{fig:response-network}



\subsubsection{Exact fractions and series coefficients}
We evaluate the tensor network in two ways to compute the LER. 
\begin{enumerate}
\item At a fixed circuit-level noise strength $p$ expressed as a fraction, we evaluate the whole contraction modulo several certified primes and reconstruct the acceptance and accepted-error probabilities as exact fractions. 
Their ratio is $\PL(p)$ with every fault order included.
Exact rational arithmetic on the whole network is impractical, the fractions run to thousands of digits, whereas modular arithmetic uses fixed-size integers and recovers the exact fraction at the end (Appendix~\ref{app:arithmetic}).
\item As a series expansion, we remove the common no-fault factor and use
polynomials in parameter $x=p/(1-p)$. We keep degrees 0 to 10 in $x$ in every factor, including the earlier-stage weights and final responses. This returns the exact raw acceptance and
accepted-error coefficients through that order. Dividing their series
gives Eq.~\eqref{eq:fullseries}. The contribution of the higher-order terms is bounded (Appendix~\ref{app:arithmetic}).
\end{enumerate}

